\ticket{10}{Продукционная модель представления знаний: структура и цикл работы продукционной системы. Нечеткие знания и их обработка в продукционных системах}

\begin{examplan}
    \item Определить продукционную систему и структуру продукционного правила.
    \item Описать базу правил, рабочую память и механизм вывода.
    \item Раскрыть цикл работы продукционной системы и направления вывода.
    \item Объяснить обработку нечетких знаний в продукционных системах.
\end{examplan}

\subsection*{Короткая версия для письменного ответа}

\textbf{Продукционная модель} представляет знания в виде правил вида \textit{если условие, то действие}. Продукционная система состоит из базы правил, рабочей памяти с текущими фактами и механизма вывода, который сопоставляет правила с фактами и применяет подходящие правила.

\subsection*{Структура продукционной системы}

\begin{description}
    \item[База правил] содержит продукции вида \texttt{IF условие THEN действие}.
    \item[Рабочая память] содержит текущие факты задачи: исходные данные и промежуточные выводы.
    \item[Механизм вывода] управляет применением правил к фактам.
\end{description}

Левая часть правила называется \textbf{антецедентом}; она задает условия применимости. Правая часть называется \textbf{консеквентом}; она может добавлять, изменять или удалять факты, запускать внешние процедуры или завершать вывод.

\subsection*{Цикл работы}

Механизм вывода работает по циклу \textbf{сопоставление -- выбор -- выполнение}.

\begin{enumerate}
    \item \textbf{Сопоставление}: условия правил сравниваются с фактами в рабочей памяти.
    \item \textbf{Формирование конфликтного множества}: находятся все правила, которые могут сработать.
    \item \textbf{Разрешение конфликтов}: выбирается одно правило или одна его инстанциация.
    \item \textbf{Выполнение}: правая часть выбранного правила изменяет рабочую память или выполняет действие.
    \item Цикл повторяется до достижения цели, пустого конфликтного множества или команды остановки.
\end{enumerate}

Стратегии разрешения конфликтов: приоритеты правил, специфичность, свежесть данных, порядок правил или эвристический выбор.

\subsection*{Прямой и обратный вывод}

\begin{description}
    \item[Прямой вывод] идет от известных фактов к новым заключениям. Подходит для мониторинга, интерпретации данных и планирования.
    \item[Обратный вывод] идет от цели к условиям, которые должны ее подтвердить. Подходит для диагностики и консультационных систем.
\end{description}

\subsection*{Достоинства и недостатки}

Достоинства:
\begin{itemize}
    \item правила просты и близки к экспертным формулировкам;
    \item базу знаний легко пополнять и изменять;
    \item возможно объяснение вывода через сработавшие правила;
    \item применимы прямой, обратный и смешанный вывод.
\end{itemize}

Недостатки:
\begin{itemize}
    \item при большом числе правил вывод может быть медленным;
    \item трудно гарантировать полноту и непротиворечивость базы правил;
    \item управление конфликтами и метаправила усложняют систему;
    \item независимые правила могут неожиданно взаимодействовать.
\end{itemize}

\subsection*{Нечеткие знания}

Классическая продукционная система работает с булевыми фактами. В реальных задачах часто встречаются нечеткие понятия: \textit{высокая температура}, \textit{молодой пациент}, \textit{средний риск}.

\textbf{Нечеткая продукционная система} использует лингвистические переменные, нечеткие термы и функции принадлежности $\mu_A(x)\in[0,1]$.

Пример правила:
\begin{verbatim}
IF temperature IS high AND pressure IS increased
THEN risk IS medium
\end{verbatim}

Этапы нечеткого вывода:
\begin{enumerate}
    \item \textbf{Фаззификация}: перевод четких входных значений в степени принадлежности термам.
    \item \textbf{Агрегирование условий}: вычисление степени истинности условий, например через минимум для AND.
    \item \textbf{Активация заключений}: перенос степени истинности условия на выходной терм.
    \item \textbf{Аккумуляция}: объединение выводов нескольких правил.
    \item \textbf{Дефаззификация}: получение четкого числового решения.
\end{enumerate}

Нечеткие продукционные системы применяются в экспертных системах, управлении, поддержке принятия решений и распознавании образов.

\begin{important}
    \item Продукционное правило имеет вид "если условие, то действие".
    \item Основные компоненты: база правил, рабочая память, механизм вывода.
    \item Цикл работы: сопоставление, конфликтное множество, выбор правила, выполнение.
    \item Прямой вывод идет от фактов к заключениям, обратный --- от цели к подтверждающим фактам.
    \item Нечеткие продукционные системы работают со степенями принадлежности, а не только с истиной и ложью.
\end{important}
